\subsection{Ordinary classification over non-Noetherian henselian unions}
\label{C:henselian-flat-unions-subsection}

\begin{lemma}[Proved: reflection under a flat local map with unchanged residue]
\label{C:flat-local-residue-reflection}
Let $A\to A'$ be a flat local map, where $A$ is
Noetherian henselian local with $2\in A^\times$,
$A'$ is any local ring, and the induced residue-field
map is an isomorphism onto a common field $\kappa$.
For finite projective rational $O_m/A$ representations,
base change preserves and reflects isomorphism,
indecomposability, self-duality and the existence of a
nonsingular invariant form of either sign.
It preserves the residue endomorphism division algebra
and its adjoint involution for an indecomposable with a
chosen form.  Isometry of two actual nonsingular
symmetric forms is also preserved and reflected.
No Noetherian assumption on $A'$ is needed.
\end{lemma}
\begin{proof}
All underlying finite projectives are free, since the
rings are local.  Equivariance of a constant matrix
is a finite linear coefficient system in the finite
free integral modules $C_d$.
Flatness therefore gives, for representations $M,N$,
\[
 \operatorname{Hom}_{O_m}(M,N)\otimes_A A'
       =\operatorname{Hom}_{O_m}(M_{A'},N_{A'}).           \tag{C.3.FU1}
\]
The Hom module on the left before tensoring is finite,
since it is a submodule of a finite free module over
the Noetherian ring $A$.
Its image in the corresponding residue matrix space
is unchanged by (C.3.FU1): every coefficient of a
linear combination over $A'$ has residue in the
same $\kappa$, and that residue lifts to $A$.
An isomorphism over $A'$ thus supplies an invertible
residue matrix in this common image.  Lifting its
coefficients to $A$ gives an equivariant matrix with
unit determinant, hence an isomorphism over $A$.
Preservation is immediate.  This argument also
applies to the linear equations for invariant
symmetric or alternating forms, so existence of
both signs is preserved and reflected.
Self-duality is the special isomorphism test with
$N=M^*$.

For $E=\operatorname{End}(M)$ and
$E'=\operatorname{End}(M_{A'})$, let $C$ be their
common image in the residue matrix algebra.
The kernels of $E\to C$ and $E'\to C$ lie in the
respective Jacobson radicals: if a matrix $h$ has
zero residue, then $1-uh$ has unit determinant
for every equivariant $u$, and its inverse is
again equivariant.  It follows that
\[
 E/\operatorname{rad}E\simeq
 C/\operatorname{rad}C\simeq
 E'/\operatorname{rad}E'.                                \tag{C.3.FU2}
\]
The source ring $E$ is semiperfect by
Lemma~\ref{C:henselian-finite-algebra-lifting}.
If $M$ is indecomposable, $E$ is local, so the
middle quotient in (C.3.FU2) is a division algebra.
The same is then true of $E'/\operatorname{rad}E'$,
which makes $E'$ local.  Hence $M_{A'}$ remains
indecomposable.  A nontrivial decomposition over
$A$ stays nontrivial after this faithfully flat
local map, proving reflection as well.
This also proves preservation of the stated
residue division algebra.  A chosen form and
its adjoint commute with base change, so the
involution in (C.3.FU2) is preserved.

Finally suppose two symmetric forms become
isometric over $A'$.  Their underlying modules
are already isomorphic over $A$, by the first
paragraph.  Use such an isomorphism to put both
forms on one module, writing them $q$ and $q a$,
where $a$ is an invertible self-adjoint element
of $E$ for the $q$-adjoint.
The isometry over $A'$ gives a residue isometry
in the common quotient (C.3.FU2).
Lemma~\ref{C:henselian-finite-algebra-lifting},
applied over the source $A$, lifts that residue
isometry to an actual isometry over $A$.
No radical lifting theorem over the possibly
non-Noetherian target was invoked.
\end{proof}

\begin{theorem}[Proved: complete ordinary algorithms for effective flat unions]
\label{C:effective-henselian-flat-unions}
Let
\[
 R_0\longrightarrow R_1\longrightarrow R_2\longrightarrow\cdots
\]
be a computably supplied sequence of solved Noetherian
henselian local rings: each stage comes with the
uniform finite-name and full classification algorithms
of Theorems~\ref{C:effective-henselian-DVR-WG},
\ref{C:finite-field-henselization-WG},
\ref{C:arithmetic-henselization-WG}, or
\ref{C:arithmetic-strict-henselization-WG}.
Assume that the transition maps are computable on
finite names, are flat and local, and induce the
identity on one fixed effective residue field
$\kappa$, either a finite field of odd characteristic
or the specified $\overline{\mathbf F}_p$.
Put $R=\operatorname*{colim}_nR_n$.

Then $R$ is henselian local and has exact finite-name
arithmetic, equality, residue and full matrix-kernel
algorithms.  No Noetherian hypothesis on $R$ is required.
All rational $O_m/R$ indecomposables have a computable
literal canonical index set, and all the isomorphism,
form-type, isometry, split $WG$ equality, multiplication
and exterior algorithms remain ordinary and terminating.
Its additive canonical description is (C.3.HD2) in
finite residue and (C.3.SH4) in algebraically closed
residue, using the resulting indices over $R$.
Moreover all transition maps and the maps to $R$ induce
injections on the split groups, and
\[
 WG_{\rm sp}(BO_{m,R})
       =\operatorname*{colim}_n WG_{\rm sp}(BO_{m,R_n})    \tag{C.3.FU3}
\]
as rings with all exterior operations.
The canonical indices and equality algorithms here are
specified explicitly; (C.3.FU3) is not being used as an
unresolved colimit-equality prescription.
\end{theorem}
\begin{proof}
Every stage map is faithfully flat, hence injective.
The colimit is local: an element represented in $R_n$
is a unit exactly when its residue in the common
$\kappa$ is nonzero.  Its maximal ideal consists of
the elements with zero residue, and its residue
field is $\kappa$.
For each $n$ the map $R_n\to R$ is flat, as a
filtered colimit of flat $R_n$-algebras, and is
local, hence faithfully flat.
The ring $R$ is henselian.  A monic polynomial
and a simple residual root use finitely many
coefficients in one $R_n$, and that root has a
lift in $R_n$ because its residue is the same
$\kappa$.  Henselianity of $R_n$ supplies the
required root there.  This proves the simple-root
criterion for $R$.

An element of $R$ is named by a stage and a
finite name in that stage.  To compare two such
elements, compute their images at the larger
stage.  Further maps are injective, so equality
in that stage is exactly equality in $R$.
Arithmetic and inverses of units are computed
at a common stage, and residues are unchanged.
Raw finite codes have decidable validity; keeping
only a code with no smaller equal predecessor
gives a decidable set of canonical names, including
when the colimit is finite.
For a finite matrix choose one stage containing
all its entries, compute kernel generators there,
and base change.  Flatness proves that they are
generators of its full actual kernel over $R$.
Thus these are ordinary algorithms, without a
search for eventual equality or stabilization.

\emph{Every object and form occurs at one stage.}
An idempotent presentation of a finite projective
module, its coaction coefficients in some finite
$C_d$, and the finitely many coefficients of a
Gram matrix and its inverse on the projective
image all occur at one stage.
Their identities hold there: the stage-to-colimit
map is injective and all coefficient modules
in the identities are finite free.
The projective inverse-Gram identities descend
with these finite data.  For a free matrix, a
determinant which is a unit in $R$ already has
nonzero common residue, and hence is a unit
at that stage as well.
Thus every actual symmetric form and every
representation descends to a stage.
Since all these local rings have free finite
projectives, free bases may be chosen in that
stage and then used over $R$.

Apply Lemma~\ref{C:flat-local-residue-reflection}
to $R_n\to R$ and to every later stage map.
It proves that isomorphism, indecomposability,
self-duality and actual form isometry are already
decided at any common stage of their data.
A Krull--Schmidt decomposition computed at a
stage remains a decomposition into indecomposables
over $R$, with local endomorphism rings by
(C.3.FU2).  The usual finite exchange argument
therefore proves uniqueness over $R$ as well:
in the expression of an identity through a
second decomposition, one summand in a local
endomorphism ring is a unit, giving a matching
indecomposable; cancel it and induct.
No global Noetherian Krull--Schmidt hypothesis
has been imposed on $R$.
The same lemma identifies all residue division
algebras and form types with the already computed
stage data.

\emph{Literal canonical indices and algorithms.}
Code finite coaction matrices by natural numbers
using the raw coefficient names just given.
Their validity and minimum coefficient degree
are decidable at a common stage.  A coefficient
outside $C_{d-1}$ cannot become zero later,
by injectivity; degree is also invariant under
constant changes of basis.
Indecomposability and isomorphism are decided by
the common-stage algorithms and the reflection
lemma.  Retain a valid indecomposable code exactly
when no smaller valid code is isomorphic to it.
This finite predecessor test is a decision
procedure, and every isomorphism class has one
least code.  It computes the claimed index set
and the index of every input module.

To choose a reference form on a self-dual index,
first construct one at its stage using the
stage algorithms.  Its finite code bounds the
search for the least valid form code of the
chosen sign; each predecessor test is again
performed at a common finite stage.
The residue division field, involution, signs
and nonsquare bits, when present, are the stage
data from (C.3.FU2).  All actual lift constructions
can be done at that stage.  If canonical choices
of references at two stages differ, base change
the actual generators and compute their coordinates
at a common stage; this is an ordinary transition
algorithm, not an assumption that arbitrary
reference choices coincide.
Consequently the already proved local residue-form
classification gives exactly (C.3.HD2) or (C.3.SH4)
with these literal indices.

Given finitely many inputs, tensor products,
exterior minors, and all specified virtual
exterior coefficients are computed at a common
stage, followed by the canonicalization just
described.  Equality and actual isometry are
decided at a common stage by the reflection
lemma.  Every step terminates.

For completeness, the monoid of actual symmetric
forms over $R$ is the filtered union of the stage
monoids: every form descends, and isometry is
reflected.  Each stage monoid is cancellative
by its residue-form classification.  Their union
is cancellative as well, since any proposed
cancellation identity occurs at a common stage.
Its group completion is therefore the filtered
union of the injective stage group completions.
This proves (C.3.FU3) and all its injectivity
claims.  Tensor and exterior constructions
commute with the stage maps, proving the stated
compatibility of operations.
\end{proof}

\begin{corollary}[Proved: infinite-dimensional local domain families]
\label{C:infinite-variable-henselian-WG}
Let $p>2$ and $q$ be a power of $p$.
Each of the following rings has the complete ordinary
canonical classification and all split $WG$ operations
of Theorem~\ref{C:effective-henselian-flat-unions}:
\[
 \begin{split}
 &\bigl(\mathbf F_q[x_1,x_2,\ldots]_{(x_1,x_2,\ldots)}\bigr)^h,
 \qquad
 \bigl(\mathbf F_q[x_1,x_2,\ldots]_{(x_1,x_2,\ldots)}\bigr)^{\rm sh},\\
 &\bigl(\mathbf Z[x_1,x_2,\ldots]_{(p,x_1,x_2,\ldots)}\bigr)^h,
 \qquad
 \bigl(\mathbf Z[x_1,x_2,\ldots]_{(p,x_1,x_2,\ldots)}\bigr)^{\rm sh}.
 \end{split}                                             \tag{C.3.FU4}
\]
The strict henselizations use the fixed algebraic
closure of the residue field.
These are countable henselian local domains of
infinite Krull dimension and infinite embedding
dimension.  In particular they are non-Noetherian,
and admit no profinite ring topology.
\end{corollary}
\begin{proof}
Let $A_n=\mathbf F_q[x_1,\ldots,x_n]_{(x_1,\ldots,x_n)}$
in the first pair, and
$A_n=\mathbf Z[x_1,\ldots,x_n]_{(p,x_1,\ldots,x_n)}$
in the second.  These regular local rings have
uniform explicit finite presentations at their
closed points.  Their local polynomial transition
maps are flat and computable.
Take $R_n=A_n^h$, or respectively $R_n=A_n^{\rm sh}$.
The preceding arithmetic henselization theorems
supply all stage algorithms uniformly in $n$.
The transition maps on pointed finite names are
computed by the inclusions of the polynomial
coordinates, with the residue embedding fixed.

The maps $R_n\to R_{n+1}$ are flat.
For henselizations this follows by identifying
$A_{n+1}^h$ with the henselization of
$A_n^h\otimes_{A_n}A_{n+1}$ at its specified
local point.  The first map is a base change
of the flat polynomial transition; localization
and henselization are flat.
The same argument with strict henselization
applies to the strict stages
\cite[Section 10.155, Lemmas 10.155.8 and 10.155.12]
 {C-StacksFlatUnions}.
They are local and induce the identity on
$\mathbf F_q$, $\mathbf F_p$, or the chosen
$\overline{\mathbf F}_p$, as appropriate.

The colimit of the $R_n$ is the corresponding
ring in (C.3.FU4).  Indeed the colimit is
henselian, or strictly henselian with the fixed
residue closure, by the simple-root argument in
the theorem.  The universal properties of the
local henselizations give maps both ways between
that colimit and the displayed henselization
of $\operatorname*{colim}A_n$.
Their composites restrict to the identity at
every finite stage by the uniqueness of the
pointed henselian lift, and hence are identities.
In the strict case all these maps induce the
identity on the chosen residue closure, so the
same uniqueness applies.  This proves the
claimed identification and all algorithmic
assertions by the theorem.

The rings $R_n$ are regular local domains, and
their cotangent spaces have bases given by
$x_1,\ldots,x_n$, or by $p,x_1,\ldots,x_n$
in mixed characteristic
\cite[Lemmas 15.46.6 and 15.46.10, Section 15.46]
 {C-StacksFlatUnions}.
All stage maps are injective, so their colimit
is a domain.  The transition maps on these
cotangent spaces are the inclusions of the
indicated bases.  The maximal ideal of the
colimit is the union of the stage maximal
ideals, and its square is the union of their
squares: every finite sum of products occurs
at a common stage.  Its cotangent space is
therefore the direct limit of these spaces
and has infinite dimension over the residue
field.  Its maximal ideal cannot be finitely
generated, proving non-Noetherianity.

Each $R_n\to R$ is faithfully flat.  Going
down from a prime over the stage maximal ideal
lifts chains of the lengths $\dim R_n$, which
are unbounded.  Thus $R$ has infinite Krull
dimension.  It is countable because it has a
countable union of finite-name presentations,
and infinite already because of its variables.
An infinite profinite additive group is
uncountable, as proved in
Theorem~\ref{C:henselian-products-WG}.
Hence these rings admit no profinite ring
topology.  This proves every assertion.
\end{proof}
